\chapter{Teilfolgengrenzwerte und Cauchy-Folgen}
\label{b21guide:ch:compactness}

Alle Aussagen dieses Kapitels betreffen reelle Folgen.

\section{Teilfolgenkriterien}

\begin{theorem}[Grenzwert jeder Teilfolge]
\label{b21guide:thm:subsequence-limit}
Ist $a_n\to L$, so gilt $a_{\varphi(n)}\to L$ für jede Teilfolge.
Besitzt eine Folge zwei Teilfolgen mit verschiedenen reellen
Grenzwerten, so ist sie divergent.
Auch eine unbeschränkte Teilfolge schließt Konvergenz aus.
\end{theorem}

\begin{theorem}[Teilteilfolgenkriterium]
\label{b21guide:thm:subsub}
Für $L\in\mathbb R$ sind äquivalent:
\begin{enumerate}
 \item $a_n\to L$.
 \item Jede Teilfolge von $(a_n)$ besitzt ihrerseits eine
 Teilfolge, die gegen $L$ konvergiert.
\end{enumerate}
Eine Beschränktheitsannahme ist für dieses Kriterium nicht nötig.
\end{theorem}

\begin{theorem}[Aufteilung in endlich viele Teilfolgen]
\label{b21guide:thm:residues}
Für jede feste ganze Zahl $r\geq1$ gilt
\[
 a_n\to L
 \quad\Longleftrightarrow\quad
 a_{rn+j}\to L\ \text{für jedes }j=0,\ldots,r-1.
\]
Insbesondere konvergiert eine Folge genau dann gegen $L$,
wenn ihre gerade und ihre ungerade Teilfolge beide gegen $L$
konvergieren.
\end{theorem}

\section{Häufungswerte}

\begin{definition}[Häufungswert einer Folge]
\label{b21guide:def:cluster}
Eine reelle Zahl $h$ heißt \emph{Häufungswert} von $(a_n)$,
wenn eine Teilfolge gegen $h$ konvergiert.
Die Menge aller reellen Häufungswerte wird mit $H(a)$ bezeichnet.
$\pm\infty$ gehören nach dieser Definition nicht zu $H(a)$.
\end{definition}

\begin{theorem}[Indexkriterium für Häufungswerte]
\label{b21guide:thm:cluster-test}
Für $h\in\mathbb R$ gilt
\[
 h\in H(a)
 \quad\Longleftrightarrow\quad
 \forall\varepsilon>0\;\forall N\in\mathbb N\;
 \exists n\geq N:\ |a_n-h|<\varepsilon.
\]
Damit muß jede solche Umgebung an unendlich vielen \emph{Indizes}
getroffen werden; die zugehörigen Werte müssen nicht verschieden sein.
\end{theorem}

\begin{theorem}[Monotone Teilfolge]
\label{b21guide:thm:monotone-subsequence}
Jede reelle Folge besitzt eine monoton wachsende oder eine
monoton fallende Teilfolge. Monotonie ist hier nicht strikt:
die Glieder der Teilfolge dürfen gleich sein.
\end{theorem}

\begin{theorem}[Erdős--Szekeres für endliche reelle Folgen]
\label{b21guide:thm:erdos-szekeres}
Seien $r,s\geq1$ ganze Zahlen. Jede endliche reelle Folge
$(a_0,\ldots,a_{N-1})$ mit $N\geq rs+1$ enthält
\begin{itemize}
 \item eine monoton wachsende Teilfolge der Länge $r+1$ oder
 \item eine monoton fallende Teilfolge der Länge $s+1$.
\end{itemize}
Eine endliche Teilfolge der Länge $\ell$ hat hier die Form
$(a_{i_0},\ldots,a_{i_{\ell-1}})$ mit
$0\leq i_0<\cdots<i_{\ell-1}<N$.
Die Monotonie ist zunächst nicht strikt; wiederholte Werte sind erlaubt.
Sind alle Folgenglieder paarweise verschieden, sind die
jeweiligen Monotonieaussagen strikt.

Die Schranke ist optimal: Es gibt bereits unter den Folgen mit
paarweise verschiedenen Gliedern solche der Länge $rs$, die
weder eine wachsende Teilfolge der Länge $r+1$ noch eine
fallende Teilfolge der Länge $s+1$ enthalten.
\end{theorem}

\begin{theorem}[Bolzano--Weierstraß]
\label{b21guide:thm:bw}
Jede beschränkte reelle Folge besitzt eine konvergente Teilfolge.
Insbesondere ist $H(a)\ne\varnothing$ für jede beschränkte
reelle Folge.

Für jede feste endliche Anzahl beschränkter reeller Folgen
$(a_n^{(1)}),\ldots,(a_n^{(m)})$, $m\geq1$, gibt es eine
gemeinsame Teilfolgenabbildung $\varphi$ und reelle Zahlen
$L_1,\ldots,L_m$ mit
\[
 a_{\varphi(n)}^{(j)}\longrightarrow L_j
 \qquad(j=1,\ldots,m).
\]
\end{theorem}

\begin{theorem}[Konvergenz durch genau einen Häufungswert]
\label{b21guide:thm:unique-cluster}
Für jede beschränkte reelle Folge und $L\in\mathbb R$ gilt
\[
 a_n\to L\quad\Longleftrightarrow\quad H(a)=\{L\}.
\]
Ohne die Beschränktheitsvoraussetzung ist die Rückrichtung
im allgemeinen falsch.
\end{theorem}

\begin{theorem}[Stabilität der Häufungswerte]
Endlich viele Änderungen von Folgengliedern ändern $H(a)$ nicht;
jede Restfolge hat dieselbe Häufungswertmenge.
Gilt $a_n-b_n\to0$, so ist $H(a)=H(b)$.
\end{theorem}

\section{Cauchy-Kriterium}

\begin{definition}[Cauchy-Folge]
\label{b21guide:def:cauchy}
Eine reelle Folge $(a_n)$ heißt \emph{Cauchy-Folge}, wenn
\[
 \forall\varepsilon>0\;\exists N\in\mathbb N\;
 \forall m,n\geq N:\ |a_m-a_n|<\varepsilon.
\]
Die Abschätzung muß für alle Paare hinreichend großer Indizes gelten.
\end{definition}

\begin{theorem}[Folgerungen aus der Cauchy-Eigenschaft]
\label{b21guide:thm:cauchy-basic}
Jede konvergente reelle Folge ist eine Cauchy-Folge.
Jede reelle Cauchy-Folge ist beschränkt, und jede ihrer Teilfolgen
ist wieder eine Cauchy-Folge. Endlich viele Änderungen und
das Entfernen endlich vieler Anfangsglieder erhalten die
Cauchy-Eigenschaft in beiden Richtungen.
\end{theorem}

\begin{theorem}[Cauchy-Folge mit konvergenter Teilfolge]
Ist $(a_n)$ eine Cauchy-Folge und konvergiert eine ihrer
Teilfolgen gegen $L$, so gilt bereits $a_n\to L$.
\end{theorem}

\begin{theorem}[Cauchy-Kriterium der reellen Zahlen]
\label{b21guide:thm:cauchy-complete}
Eine reelle Folge ist genau dann konvergent, wenn sie eine Cauchy-Folge
ist. Die Existenz eines reellen Grenzwerts verwendet die Vollständigkeit
der reellen Zahlen.
\end{theorem}

\begin{theorem}[Benachbarte Differenzen]
Für jede reelle Cauchy-Folge gilt
$a_{n+1}-a_n\to0$.
Diese Bedingung allein ist nicht hinreichend für die
Cauchy-Eigenschaft oder für Konvergenz.
\end{theorem}

\section{Weitere Häufungswert- und Konvergenzkriterien}

\begin{theorem}[Endlicher Wertebereich und Periodizität]
\label{b21guide:thm:finite-range}
Sei die Wertemenge $\{a_n:n\in\mathbb N\}$ endlich.
Dann besteht $H(a)$ genau aus den Werten, die unendlich oft
vorkommen, und
\[
 (a_n)\text{ konvergent}
 \quad\Longleftrightarrow\quad
 (a_n)\text{ schließlich konstant}.
\]
Insbesondere ist eine schließlich periodische reelle Folge
genau dann konvergent, wenn sie schließlich konstant ist.
Eine von Anfang an periodische reelle Folge ist genau dann
konvergent, wenn sie konstant ist.
\end{theorem}

\begin{theorem}[Grenzwerte von Häufungswerten]
\label{b21guide:thm:closed-cluster}
Sei $(h_j)$ eine reelle Folge mit $h_j\in H(a)$ für alle $j$.
Aus $h_j\to h\in\mathbb R$ folgt $h\in H(a)$.
Dieser Satz setzt keine Beschränktheit von $(a_n)$ voraus.
\end{theorem}

\begin{theorem}[Zwischenwerte bei verschwindenden Schritten]
\label{b21guide:thm:cluster-intermediate}
Gelte $a_{n+1}-a_n\to0$.
Sind $u,v\in H(a)$ und $u<v$, so ist jede reelle Zahl $t$
mit $u\leq t\leq v$ ebenfalls in $H(a)$.
Bei einer beschränkten Folge mit verschwindenden benachbarten
Differenzen besteht die Häufungswertmenge somit entweder aus
genau einem Wert oder enthält ein ganzes nichttriviales Intervall.
\end{theorem}

\begin{theorem}[Geometrisch abnehmende Schritte und Fehlerabschätzung]
\label{b21guide:thm:geometric-error}
Seien $C\geq0$, $0<q<1$ und $r\in\mathbb N$ gegeben.
Gelte für alle $n\geq r$
\[
 |a_{n+1}-a_n|\leq Cq^{\,n-r}.
\]
Dann konvergiert $(a_n)$ gegen ein $L\in\mathbb R$, und für
$n\geq r$ gilt die Fehlerabschätzung
\[
 |a_n-L|\leq\frac{Cq^{\,n-r}}{1-q}.
\]
Insbesondere genügt
$|a_{n+2}-a_{n+1}|\leq q\,|a_{n+1}-a_n|$ für alle $n\geq r$;
dann kann $C=|a_{r+1}-a_r|$ gewählt werden, und zusätzlich gilt
\[
 |a_n-L|\leq\frac{q}{1-q}|a_n-a_{n-1}|
 \qquad(n\geq r+1).
\]
\end{theorem}

\clearpage
\section{Indexlücken und stabile Rekursionen}

\begin{theorem}[Konvergenz über eine Teilfolge mit beschränkten Lücken]
\label{b21guide:thm:bounded-gaps}
Es gelte $a_{n+1}-a_n\to0$.
Sei $(n_j)$ eine strikt wachsende Folge natürlicher Zahlen, und
es gebe eine ganze Zahl $p\geq1$ mit
$n_{j+1}-n_j\leq p$ für alle hinreichend großen $j$.
Dann gilt für jedes $L\in\mathbb R$
\[
 a_{n_j}\to L\quad\Longrightarrow\quad a_n\to L.
\]
Dasselbe gilt für bestimmte Divergenz:
$a_{n_j}\to+\infty$ erzwingt $a_n\to+\infty$,
und $a_{n_j}\to-\infty$ erzwingt $a_n\to-\infty$.
Eine Beschränktheit der ganzen Folge wird nicht vorausgesetzt.
\end{theorem}

\begin{theorem}[Stabile Fehlerrekursion und affine Iteration]
\label{b21guide:thm:perturbed-recursion}
Seien $L\in\mathbb R$, $0<q<1$, $r\in\mathbb N$ und
$\varepsilon_n\geq0$ mit $\varepsilon_n\to0$.
Gelte für alle $n\geq r$
\[
 |a_{n+1}-L|\leq q|a_n-L|+\varepsilon_n.
\]
Dann ist $a_n\to L$. Für $n\geq r$ gilt die endliche Fehlerabschätzung
\[
 |a_n-L|
 \leq q^{\,n-r}|a_r-L|
       +\sum_{k=r}^{n-1}q^{\,n-1-k}\varepsilon_k,
\]
wobei die Summe für $n=r$ gleich $0$ ist.

Insbesondere sei eine reelle Folge durch
\[
 a_{n+1}=c_na_n+d_n,\qquad a_0\in\mathbb R,
\]
gegeben. Aus $c_n\to c$ mit $|c|<1$ und $d_n\to d\in\mathbb R$
folgt für jeden Anfangswert
\[
 a_n\longrightarrow\frac{d}{1-c}.
\]
Die Koeffizienten dürfen vom Index abhängen.
Die Bedingung $|c|<1$ ist für dieses allgemeine Kriterium wesentlich.
\end{theorem}
